% from ega2-1-fr.tex
\oldpage[II]{66}

que nous appellerons le morphisme \emph{associé à $\mathcal{L}$ et à $\psi$}~;
rappelons (\hyperref[II.3.5.6-fr]{3.5.6}) que ce morphisme s'obtient en
composant avec la première projection
$\pi:\operatorname{Proj}(q^*(\mathcal{S}))=P\times_YX\to P$ le
$Y$-morphisme
\[
  \tau=\operatorname{Proj}(\psi):G(\psi)\longrightarrow
  \operatorname{Proj}(q^*(\mathcal{S})).
\]
\end{env}

\begin{env}[3.7.2]
\label{II.3.7.2-fr}
Explicitons $r=r_{\mathcal{L},\psi}$ lorsque $Y=\operatorname{Spec}(A)$ est
affine, et par suite $\mathcal{S}=\widetilde{S}$, où $S$ est une $A$-algèbre
graduée à degrés positifs. Supposons d'abord que $X=\operatorname{Spec}(B)$
soit aussi affine et que l'on ait $\mathcal{L}=\widetilde{L}$, où $L$ est un
$B$-module libre de rang $1$. On a alors
$q^*(\mathcal{S})=(S\otimes_AB)^{\sim}$
(\hyperref[I.1.6.5-fr]{I, 1.6.5})~; si $c$ est un générateur de $L$,
$\psi_n:q^*(\mathcal{S}_n)\to\mathcal{L}^{\otimes n}$ correspond à un
homomorphisme
$w_n:s\otimes b\mapsto bv_n(s)c^{\otimes n}$ de $S_n\otimes_AB$ dans
$L^{\otimes n}$, où $v_n:S_n\to B$ est un homomorphisme de $A$-modules, les
$v_n$ constituant un homomorphisme d'algèbres $S\to B$. Soit $f\in S_d$
($d>0$) et posons $g=v_d(f)$~; on a
$\pi^{-1}(D_+(f))=D_+(f\otimes1)$ en vertu de
(\hyperref[II.2.8.10-fr]{2.8.10}) et de l'identification de $D_+(f)$ à
$\operatorname{Spec}(S_{(f)})$ (\hyperref[II.2.3.6-fr]{2.3.6})~; d'autre
part, la formule (\hyperref[II.2.8.1.1-fr]{2.8.1.1}) montre (compte tenu de
l'identification canonique de $X$ et de $\operatorname{Proj}(\mathcal{S}')$)
que
\[
  \tau^{-1}(D_+(f\otimes1))=D(g)
\]
d'où
\[
\label{II.3.7.2.1-fr}
  r^{-1}(D_+(f))=D(g).
\tag{3.7.2.1}
\]

En outre, le morphisme $\tau=\operatorname{Proj}(\psi)$, restreint à $D(g)$,
correspond à l'homomorphisme qui applique
$(s\otimes1)/(f\otimes1)^n$ (pour $s\in S_{nd}$) sur
$v_{nd}(s)/g^n$ (\hyperref[II.2.8.1-fr]{2.8.1}), et la projection $\pi$,
restreinte à $D_+(f\otimes1)$, correspond à l'homomorphisme
$s/f^n\mapsto(s\otimes1)/(f\otimes1)^n$~; on en conclut que $r$, restreint à
$D(g)$, correspond à l'homomorphisme de $A$-algèbres
$\omega:S_{(f)}\to B_g$ tel que $\omega(s/f^n)=v_{nd}(s)/g^n$ pour
$s\in S_{nd}$ ($n>0$). Passant au cas où $X$ est quelconque ($Y$ étant
toujours affine), on obtient donc, compte tenu de
(\hyperref[II.2.8.1-fr]{2.8.1})~:
\end{env}

\begin{proposition}[3.7.3]
\label{II.3.7.3-fr}
Si $Y=\operatorname{Spec}(A)$ est affine et $\mathcal{S}=\widetilde{S}$, où
$S$ est une $A$-algèbre graduée, pour tout
$f\in S_d=\Gamma(Y,\mathcal{S}_d)$, on a
\[
\label{II.3.7.3.1-fr}
  r_{\mathcal{L},\psi}^{-1}(D_+(f))=X_{\psi^\flat(f)}
  \qquad\bigl(\text{où }\psi^\flat(f)\in
  \Gamma(X,\mathcal{L}^{\otimes d})\bigr)
\tag{3.7.3.1}
\]
et la restriction $X_{\psi^\flat(f)}\to D_+(f)=\operatorname{Spec}(S_{(f)})$
de $r_{\mathcal{L},\psi}$ correspond
(\hyperref[I.2.2.4-fr]{I, 2.2.4}) à l'homomorphisme d'algèbres
\[
\label{II.3.7.3.2-fr}
  \psi_{(f)}^\flat:S_{(f)}\longrightarrow
  \Gamma(X_{\psi^\flat(f)},\mathcal{O}_X)
\tag{3.7.3.2}
\]
tel que, pour $s\in S_{nd}=\Gamma(Y,\mathcal{S}_{nd})$,
\[
\label{II.3.7.3.3-fr}
  \psi_{(f)}^\flat(s/f^n)
  =(\psi^\flat(s)|X_{\psi^\flat(f)})
   (\psi^\flat(f)|X_{\psi^\flat(f)})^{-n}.
\tag{3.7.3.3}
\]
\end{proposition}

Nous dirons que $r_{\mathcal{L},\psi}$ est \emph{partout défini} si l'on a
$G(\psi)=X$. Pour qu'il en soit ainsi, il faut et il suffit évidemment que
$G(\psi)\cap q^{-1}(U)=q^{-1}(U)$ pour tout ouvert affine $U\subset Y$~;
autrement dit, la question est \emph{locale sur $Y$}. Si $Y$ est affine,
$G(\psi)$ est réunion des $r^{-1}(D_+(f))$ pour $f$ homogène dans $S_+$
(\hyperref[II.2.8.1-fr]{2.8.1})~; en vertu de
(\hyperref[II.3.7.3.1-fr]{3.7.3.1}), les $X_{\psi^\flat(f)}$ doivent donc
former un \emph{recouvrement} de $X$, autrement dit~:

\begin{corollary}[3.7.4]
\label{II.3.7.4-fr}
Sous les hypothèses de (\hyperref[II.3.7.3-fr]{3.7.3}), pour que
$r_{\mathcal{L},\psi}$ soit partout défini, il faut et il suffit que pour tout
$x\in X$, il existe un entier $n>0$ et une section $s$ de $\mathcal{S}_n$
au-dessus de $Y$ telle que si on pose
$t=\psi^\flat(s)\in\Gamma(X,\mathcal{L}^{\otimes n})$, on ait $t(x)\neq0$.
\end{corollary}

